% Prostudováno: ucebnice/02_Funkce_a_rovnice/Linearni_funkce_a_rovnice/Linerání rovnice.tex,
% část Soustavy lineárních rovnic, řešení a dosazovací metoda (ř. 998–1123 a 1219–1337).
% Realisticky.cz: externi_zdroje/realisticky/02_Funkce_a_rovnice/03_Nektere_rovnice_a_nerovnice_prevoditelne_na_linearni/020310_Soustavy_dvou_linearnich_rovnic_o_dvou_neznamych_I/020310_Lekce.pdf
% a 020311_Soustavy_dvou_linearnich_rovnic_o_dvou_neznamych_II/020311_Lekce.pdf; využito odvození dosazení a upozornění na chybu dosazení do téže rovnice. Úlohy vlastní.
\documentclass[11pt,a4paper]{article}
% Spolecny styl pro ucitelske pripravy MAT (Work/MAT)
% Verze 1.1, 2026-10-23
\NeedsTeXFormat{LaTeX2e}
\ProvidesPackage{priprava}[2026/10/23 v1.1 Sablona ucitelskych priprav MAT]

% Zakladni typografie a cestina
\RequirePackage[T1]{fontenc}
\RequirePackage[utf8]{inputenc}
\RequirePackage[czech]{babel}
\RequirePackage{lmodern}
\RequirePackage{microtype}
\RequirePackage{amsmath,amssymb}
\RequirePackage{xcolor}
\RequirePackage{enumitem}
\RequirePackage[most]{tcolorbox}
\RequirePackage{tabularx,booktabs,array}
\RequirePackage{fancyhdr}
\RequirePackage{needspace}
\RequirePackage{geometry}
\RequirePackage[hidelinks]{hyperref}

\geometry{a4paper,top=17mm,bottom=18mm,left=23mm,right=23mm}
\definecolor{notegray}{gray}{0.39}
\setlength{\parindent}{0pt}
\setlength{\parskip}{2.5pt plus 1pt minus 0.5pt}
\setlist[itemize]{topsep=2pt,itemsep=1pt,parsep=0pt,leftmargin=1.45em}
\setlist[enumerate]{topsep=2pt,itemsep=2pt,parsep=0pt,leftmargin=1.65em}
\setcounter{secnumdepth}{0}
\renewcommand{\arraystretch}{1.12}

% Zahlavi: vlevo rocnik | globalni tema, vpravo cislo: nazev hodiny
% Zapati zustava zcela prazdne.
\pagestyle{fancy}
\fancyhf{}
\renewcommand{\headrulewidth}{0.35pt}
\renewcommand{\footrulewidth}{0pt}
\setlength{\headheight}{14pt}
\newcommand{\nastavPripravu}[4]{%
  \fancyhead[L]{\small #1\enspace\textbar\enspace #2}%
  \fancyhead[R]{\small #3: #4}%
}

% Minimalisticke cernobile boxy. Kazda definice zvlast.
\tcbset{commonlesson/.style={enhanced,breakable=false,colback=white,colframe=black,
  colbacktitle=white,coltitle=black,boxrule=0.45pt,
  arc=0.8mm,outer arc=0.8mm,left=3mm,right=3mm,top=1.6mm,bottom=1.6mm,
  toptitle=1mm,bottomtitle=0.7mm,fonttitle=\bfseries,
  attach boxed title to top left={xshift=2.5mm,yshift=-1.15mm},
  boxed title style={colback=white,colframe=white,boxrule=0pt,left=1mm,right=1mm}}}
\newtcolorbox{zadaniBox}[1][]{commonlesson,title={Zadání},#1}
\newtcolorbox{prikladBox}[1][]{commonlesson,title={Příklad},#1}
\newtcolorbox{definiceBox}[1][]{commonlesson,title={Definice},#1}

% Spolecny box pro zadani a jeho reseni, bez mezery mezi castmi.
% Pouziti:
%   \begin{prikladReseniBox}
%     Zadani ulohy.
%     \tcblower
%     \begin{align*} ... \end{align*}
%   \end{prikladReseniBox}
% Spodni cast ma jemny sedy podklad a delici linku; lze v ni sazet
% vice radku matematiky. Obe casti maji jen jeden spolecny ramecek.
\tcbset{solutionlesson/.style={
  commonlesson,
  enhanced jigsaw,
  breakable,
  colbacklower=black!3,
  segmentation style={solid,draw=black!25,line width=0.35pt},
  fontlower=\small\color{notegray},
  before lower={\noindent\textbf{Řešení:}\enspace},
  leftlower=3mm,rightlower=3mm
}}
\newtcolorbox{prikladReseniBox}[1][]{solutionlesson,title={Příklad},#1}
\newtcolorbox{zadaniReseniBox}[1][]{solutionlesson,title={Zadání},#1}

% Pomucky pro autora pripravy: sedive poznamky; \reseni pro starsi pripravy
\newcommand{\poznamka}[1]{\par\begingroup\small\color{notegray}\itshape
  \textnormal{\textbf{Poznámka pro učitele.}}\enspace #1\par\endgroup}
\newcommand{\reseni}[1]{\par\begingroup\small\color{notegray}
  \textbf{Očekávané řešení:}\enspace #1\par\endgroup}
\newcommand{\blok}[1]{\needspace{5\baselineskip}\subsection*{#1}}
\newcommand{\vyrok}[1]{\uv{#1}}

\nastavPripravu{1. ročník}{Soustavy rovnic}{77}{Dosazovací metoda}
\begin{document}
\blok{Dvě podmínky pro jednu dvojici}
Jedna rovnice se dvěma neznámými připouští zpravidla mnoho dvojic. Soustava požaduje, aby stejná dvojice splnila obě podmínky současně. Z hodiny 68 víme, že graficky jde o společný bod dvou přímek; dnes jej najdeme přesně dosazováním.
\begin{definiceBox}
\textbf{Soustava dvou lineárních rovnic o dvou neznámých} je dvojice rovnic $a_1x+b_1y=c_1$ a $a_2x+b_2y=c_2$, kde $a_i,b_i,c_i\in\mathbb R$ a v každé rovnici je alespoň jeden z koeficientů u neznámých nenulový.
\end{definiceBox}
\begin{definiceBox}
\textbf{Řešení soustavy} je uspořádaná dvojice $(x_0;y_0)$, která splňuje \emph{obě} rovnice. Množina všech řešení je $K$; u jediného řešení píšeme například $K=\{(4;3)\}$.
\end{definiceBox}
\begin{definiceBox}
\textbf{Dosazovací metoda}: z jedné rovnice vyjádříme jednu neznámou, tento výraz dosadíme za tutéž neznámou do \emph{druhé} rovnice, dopočteme obě hodnoty a ověříme je v původní soustavě.
\end{definiceBox}
\begin{prikladReseniBox}
Řešte dosazováním soustavu $x+y=7$, $x-y=1$. Proveďte zkoušku v obou rovnicích.
\tcblower
Z první rovnice $x=7-y$. Dosadíme do druhé: $(7-y)-y=1$, tedy $7-2y=1$, $y=3$. Zpět $x=7-3=4$. Zkouška: $4+3=7$ a $4-3=1$. Tedy $K=\{(4;3)\}$.
\end{prikladReseniBox}
\begin{prikladReseniBox}
Řešte $2x+y=8$, $x-2y=-1$. Kterou neznámou je výhodné vyjádřit?
\tcblower
Z první rovnice je snadné vyjádřit $y=8-2x$. Dosazení do druhé dává $x-2(8-2x)=-1$, tedy $x-16+4x=-1$, $5x=15$, $x=3$. Pak $y=8-6=2$. Zkouška: $2\cdot3+2=8$ a $3-2\cdot2=-1$; $K=\{(3;2)\}$.
\end{prikladReseniBox}
\poznamka{Trvej na dosazení do druhé rovnice. Dosazení vyjádřeného $x$ zpět do téže rovnice vede jen k identitě, aniž by využilo druhou podmínku.}
\blok{Samostatné procvičení}
\begin{zadaniReseniBox}
\begin{enumerate}[label=\alph*)]
\item Řešte $x+y=5$, $2x-y=1$.
\item Řešte $y=3x-2$, $x+y=6$.
\item Pro soustavu $x=2y+1$, $3x-y=8$ určete, zda je dvojice $(3;1)$ řešením, a potom najděte celé řešení.
\end{enumerate}
\tcblower
\begin{enumerate}[label=\alph*)]
\item $y=5-x$; $2x-(5-x)=1$ dává $3x=6$, $x=2$, $y=3$. Zkouška $2+3=5$, $4-3=1$.
\item $(3x-2)+x=6$, tedy $4x=8$, $x=2$, $y=4$. Kontrola $4=3\cdot2-2$ a $2+4=6$.
\item Dvojice $(3;1)$ splňuje obě rovnice: $3=2\cdot1+1$ a $3\cdot3-1=8$. Dosazení $x=2y+1$ do druhé dává $3(2y+1)-y=8$, $5y=5$, $y=1$, $x=3$. Jediné řešení je $K=\{(3;1)\}$.
\end{enumerate}
\end{zadaniReseniBox}
\blok{Volitelné rozšíření a rezerva}
\begin{zadaniReseniBox}
Řešte $3x-2y=4$, $x+y=3$ dvěma různými volbami vyjádřené neznámé.
\tcblower
Z druhé rovnice $y=3-x$: $3x-2(3-x)=4$, tedy $5x=10$, $x=2$, $y=1$. Lze také vyjádřit $x=3-y$: $3(3-y)-2y=4$, tedy $9-5y=4$, $y=1$, $x=2$. Kontrola $6-2=4$ a $2+1=3$.
\end{zadaniReseniBox}
\begin{zadaniReseniBox}
Najděte chybu v postupu: z $x+y=7$ vyjádříme $x=7-y$ a dosadíme zpět do $x+y=7$, čímž vyjde $7=7$; závěr „soustava $x+y=7$, $x-y=1$ má nekonečně mnoho řešení“.
\tcblower
Dosazovalo se do téže rovnice, z níž se $x$ vyjádřilo. Rovnice $x-y=1$ nebyla vůbec použita. Po dosazení do ní vyjde $7-2y=1$, tedy jediná dvojice $(4;3)$.
\end{zadaniReseniBox}
\blok{Domácí práce}
\begin{zadaniReseniBox}
\begin{enumerate}[label=\alph*)]
\item Řešte $x+2y=7$, $3x-y=0$.
\item Řešte $2x-y=5$, $x+y=4$.
\item Vysvětlete, proč samotná dvojice splňující první rovnici soustavy nemusí být řešením soustavy.
\end{enumerate}
\tcblower
\begin{enumerate}[label=\alph*)]
\item Z druhé rovnice $y=3x$. Dosazení: $x+6x=7$, tedy $x=1$, $y=3$. Zkouška $1+6=7$, $3-3=0$.
\item Z druhé rovnice $y=4-x$. Pak $2x-(4-x)=5$, $3x=9$, $x=3$, $y=1$. Zkouška $6-1=5$, $3+1=4$.
\item Např. $(0;7)$ splňuje $x+y=7$, ale nesplňuje $x-y=1$, protože $0-7=-7$. Soustava znamená průnik podmínek.
\end{enumerate}
\end{zadaniReseniBox}
\poznamka{Jádro: význam dvojice a metoda 11 min, dva řešené příklady 13 min, samostatné řešení a kontrola 18 min, organizace 3 min. Rezerva je nad rámec; domácí práce 15–20 min. Hodina 78 zavede jiný způsob odstranění neznámé.}
\end{document}
